% ============================================================
%  Preamble: packages, theorem environments, and macros.
% ============================================================

% --- Layout ---
\usepackage[a4paper,margin=1in]{geometry}

% --- Japanese support (requires lualatex; uses bundled Haranoaji fonts) ---
\usepackage{luatexja}

% --- Math ---
\usepackage{amsmath}
\usepackage{amssymb}
\usepackage{amsthm}
\usepackage{mathtools}

% --- Diagrams ---
\usepackage{tikz}
\usepackage{tikz-cd}

% --- Lists ---
\usepackage{enumitem}

% --- Hyperlinks (load late) ---
\usepackage[colorlinks=true,linkcolor=blue,citecolor=blue,urlcolor=blue]{hyperref}

% --- Theorem environments (shared counter, numbered per chapter) ---
\theoremstyle{plain}
\newtheorem{theorem}{Theorem}[chapter]
\newtheorem{lemma}[theorem]{Lemma}
\newtheorem{proposition}[theorem]{Proposition}
\newtheorem{corollary}[theorem]{Corollary}

\theoremstyle{definition}
\newtheorem{definition}[theorem]{Definition}
\newtheorem{example}[theorem]{Example}

\theoremstyle{remark}
\newtheorem{remark}[theorem]{Remark}

% --- Math macros ---
% hyperref's PU encoding defines \C as a bookmark accent we never use;
% reclaim it for the complex numbers (hence \renewcommand, not \newcommand).
\renewcommand{\C}{\mathbb{C}}
\newcommand{\R}{\mathbb{R}}
\newcommand{\Z}{\mathbb{Z}}
\newcommand{\N}{\mathbb{N}}
\newcommand{\dbar}{\bar{\partial}}      % Dolbeault operator
\newcommand{\Oh}{\mathcal{O}}           % structure sheaf
\DeclareMathOperator{\Hom}{Hom}
% Sheaf cohomology: \coh{q}{X}{\Oh} -> H^q(X, O)
\newcommand{\coh}[3]{H^{#1}\!\left(#2,\,#3\right)}
